% 02-related
\subsection*{Consumer and low-cost EEG validation}
Comparisons of medical-grade and consumer wireless systems report that consumer
hardware can recover common event-related components under favourable conditions,
while signal quality and artifact susceptibility differ substantially between devices
and sessions \cite{ratti2017comparison,krigolson2017muse,sawangjai2020consumer}. One of
the closest prior studies compared four systems (including MindWave and Muse) on five
participants across two visits, focusing on Fp1 and explicitly comparing eyes-closed
and eyes-open spectra and their test--retest reliability
\cite{ratti2017comparison}. That is close to our alpha arm on real hardware, so the gap
cannot be "nobody has looked at consumer spectral reliability"; the gap is that
existing work does not quantify which analyses survive a \emph{controlled} degradation
chain, and does not connect that chain to claim-level verdicts. Work on channel subsets for emotion decoding has shown that small, well-chosen
montages can be competitive with large ones \cite{zheng2015channels}; what is missing is
a controlled degradation of the same recordings with a site-matched reference and a
claim-level decision rule. Our contribution is therefore the decomposition and the
decision matrix, not the observation that fewer channels can hurt.

\subsection*{Reproducibility and reporting standards}
Psychology and neuroscience face a well-documented reproducibility problem
\cite{osc2015reproducibility}. In electrophysiology the community response has been
standardised reporting (COBIDAS MEEG) \cite{pernet2020cobidas} and standardised data
organisation (EEG-BIDS) \cite{pernet2019bids}. Both are written primarily for
research-grade acquisitions; they specify what to report, but not what can be
inferred when the acquisition is a single channel. The present benchmark is designed
to fill that inference gap and to make the reporting checklist operational for
consumer-device studies.

\subsection*{Analysis pipelines and artifact handling}
Modern open pipelines make preprocessing explicit and reproducible
\cite{gramfort2013mne,delorme2004eeglab,oostenveld2011fieldtrip}, and automated
artifact rejection and independent-component classification reduce researcher degrees
of freedom \cite{jas2017autoreject,piontonachini2019iclabel}. These tools assume
enough channels for spatial decomposition; with one channel, ICA is unavailable and
most artifact rejection reduces to amplitude thresholds or regression against
auxiliary signals. A benchmark of single-channel conditions therefore must treat
artifact handling as part of the independent variable, not as a preprocessing detail.

\subsection*{Feature validity and open data}
Band-power features are the workhorse of consumer-EEG claims, and their test--retest
reliability is a prerequisite for any individual-level interpretation. Open datasets
with repeated conditions and affective ratings allow reliability and decoding to be
measured against a known ground truth rather than asserted
\cite{schalk2009eegbci,koelstra2012deap}. We build directly on the EEGBCI resources and
on an open affective dataset that contains both low-cost and high-end recordings
\cite{sanchezreolid2022emotion,sanchezreolid2025dataset}; DEAP and DREAMER are discussed
as alternatives that could not be used. We report every degradation step so that a
reader can map our results onto their own acquisition.

\subsection*{Positioning}
Relative to device-comparison work, our contribution is analysis-centred and
degradation-based. Relative to reporting standards, it is quantitative: rather than
saying which variables to report, it estimates what each variable is worth under
single-channel constraints. To our knowledge, no existing benchmark reports, for the
same open datasets, how alpha reactivity, single-trial decoding, reliability and blink
susceptibility change as a function of explicit consumer-hardware degradations.

\subsection*{Benchmarking methodology}
The idea of comparing analysis pipelines on shared datasets is established in BCI
research: MOABB provides trustworthy benchmarking across datasets and pipelines
\cite{jayaram2018moabb}. Our benchmark differs in object --- it holds the pipeline
approximately fixed and varies the acquisition conditions, rather than comparing
algorithms --- but it inherits the principle that a claim about a method should be
evaluated against a common, open reference. Filter design is a concrete example of
why this matters: apparently small high-pass choices can create artifactual effects
and incorrect conclusions \cite{tanner2015filters}, which is exactly the kind of
degradation our chain isolates and reports.
